\chapter{Adapter Contracts}

Gert v2 decouples the core runtime from presentation adapters (TUI, Web, VS
Code) and integration systems (CI/CD platforms, observability backends, approval
workflow systems) through explicitly versioned contracts.  An \textbf{adapter
contract} is a stable, versioned interface that gert commits to maintain across
minor and patch releases.  Adapters are consumers of these contracts: they
observe, render, or orchestrate gert runs, but they own no execution logic.

This chapter defines the four primary adapter surfaces exposed by gert v2:
JSON-RPC execution contract, WebSocket event stream contract, tool invocation
contract, and extension handshake contract.

% ─────────────────────────────────────────────────────────────────────────────
\section{What is an Adapter Contract?}
\label{sec:adapter-contract-definition}
% ─────────────────────────────────────────────────────────────────────────────

An \textbf{adapter contract} is a versioned API specification that defines:

\begin{enumerate}
  \item The set of RPC methods, request/response schemas, and error envelopes.
  \item The stability guarantee: which fields are required, which are optional,
        and which may be added in future minor versions.
  \item The versioning scheme: how breaking changes are communicated, and how
        long deprecated versions are supported.
  \item The transport layer: HTTP, WebSocket, stdio JSON-RPC, etc.
\end{enumerate}

Adapter contracts are the integration surface of gert.  They are designed for
external systems to consume.  Breaking changes to a contract require a major
version bump and a deprecation period.

\paragraph{Example consumers:}

\begin{itemize}
  \item VS Code extension (consumes JSON-RPC execution contract and WebSocket
        event stream)~\cite{vscode-extension-api}
  \item GitLab CI runner (consumes JSON-RPC execution contract to trigger runs)
  \item Datadog integration (consumes event stream to emit metrics)
  \item Slack approval bot (consumes event stream, submits approvals via
        JSON-RPC)
\end{itemize}

% ─────────────────────────────────────────────────────────────────────────────
\section{Contract Categories}
\label{sec:contract-categories}
% ─────────────────────────────────────────────────────────────────────────────

Gert v2 exposes four distinct adapter surfaces, each with its own versioning
and stability guarantees.

\subsection{JSON-RPC Execution Contract (\texttt{exec/v2})}
\label{sec:jsonrpc-exec-contract}

The execution contract is the primary interface for starting, monitoring, and
controlling gert runs.  It is served over \texttt{gert serve --stdio} (stdio
JSON-RPC) or \texttt{gert serve --http} (HTTP POST to \texttt{/rpc}).

\paragraph{Contract version:} \texttt{exec/v2}

\paragraph{Transport:} stdio JSON-RPC 2.0 or HTTP POST (JSON-RPC over HTTP)

\paragraph{Methods:}

\begin{table}[ht]
\centering
\begin{tabular}{p{0.25\linewidth}p{0.60\linewidth}}
\hline
\textbf{Method} & \textbf{Description} \\
\hline
\texttt{exec/start} &
Start a new runbook execution.  Params: \texttt{runbookPath}, \texttt{inputs},
\texttt{mode}, \texttt{actorId}.  Returns: \texttt{runId}, \texttt{status}. \\
\hline
\texttt{exec/next} &
Advance the run by one step or until it reaches a pause point (approval,
manual, error).  Params: \texttt{runId}.  Returns: \texttt{stepId},
\texttt{status}, \texttt{output}. \\
\hline
\texttt{exec/cancel} &
Cancel an in-progress run.  Params: \texttt{runId}, \texttt{reason}.  Returns:
\texttt{cancelled: bool}. \\
\hline
\texttt{exec/status} &
Get the current status of a run.  Params: \texttt{runId}.  Returns:
\texttt{status}, \texttt{currentStepId}, \texttt{progress}. \\
\hline
\texttt{run/list} &
List all runs in the workspace.  Params: \texttt{filter}, \texttt{limit}.
Returns: \texttt{runs: [RunSummary]}. \\
\hline
\texttt{run/get} &
Get full details of a run.  Params: \texttt{runId}.  Returns: \texttt{run:
RunDetails}. \\
\hline
\end{tabular}
\caption{JSON-RPC execution contract methods (\texttt{exec/v2}).}
\label{tab:exec-methods}
\end{table}

\paragraph{Request schema — \texttt{exec/start}:}

\begin{minted}{json}
{
  "jsonrpc": "2.0",
  "id": "1",
  "method": "exec/start",
  "params": {
    "runbookPath": "./runbooks/incident-triage.yaml",
    "inputs": {
      "incident_id": "INC-12345",
      "severity": "high"
    },
    "mode": "real",        // real | dry-run | replay
    "actorId": "ops@acme.example"
  }
}
\end{minted}

\paragraph{Response schema — \texttt{exec/start}:}

\begin{minted}{json}
{
  "jsonrpc": "2.0",
  "id": "1",
  "result": {
    "runId": "run-abc123",
    "status": "running",
    "createdAt": "2026-04-18T12:00:00Z",
    "traceFile": "./traces/run-abc123.jsonl"
  }
}
\end{minted}

\paragraph{Error envelope:}

All errors follow JSON-RPC 2.0 error format:

\begin{minted}{json}
{
  "jsonrpc": "2.0",
  "id": "1",
  "error": {
    "code": 1001,
    "message": "run not found",
    "data": {
      "runId": "run-xyz789",
      "category": "client_error",
      "retryable": false
    }
  }
}
\end{minted}

See Section~\ref{sec:error-codes} for the canonical error code registry.

\paragraph{Stability guarantee:}

Within the \texttt{exec/v2} contract:

\begin{itemize}
  \item All required fields in request/response schemas are stable.
  \item Optional fields may be added in minor versions (clients MUST ignore
        unknown fields).
  \item Required fields will not be removed or renamed (breaking change →
        \texttt{exec/v3}).
  \item Error codes are stable; new codes may be added but existing codes will
        not change semantics.
\end{itemize}

\subsection{WebSocket Event Stream Contract (\texttt{events/v2})}
\label{sec:ws-event-contract}

The event stream contract allows adapters to subscribe to real-time events from
a running or completed gert execution.  It is served over \texttt{gert serve
--http} at the \texttt{/ws} or \texttt{/events} endpoint.

\paragraph{Contract version:} \texttt{events/v2}

\paragraph{Transport:} WebSocket (wss:// or ws://)

\paragraph{Connection lifecycle:}

\begin{enumerate}
  \item Client establishes WebSocket connection to \texttt{/ws}.
  \item Client sends a \texttt{subscribe} message with \texttt{runId}.
  \item Server streams all events for that run in order.
  \item Client may disconnect at any time; server closes gracefully.
\end{enumerate}

\paragraph{Subscribe message:}

\begin{minted}{json}
{
  "action": "subscribe",
  "runId": "run-abc123",
  "sinceSequence": 0  // Optional: replay from sequence number
}
\end{minted}

\paragraph{Event envelope:}

Events are transmitted as JSON objects, one per WebSocket message.  The event
envelope is identical to the trace event format defined in §12:

\begin{minted}{json}
{
  "eventType": "step/started",
  "timestamp": "2026-04-18T12:00:00Z",
  "runId": "run-abc123",
  "stepId": "lookup_incident",
  "sequence": 42,
  "payload": {
    "stepType": "tool",
    "toolName": "acme.pd-lookup"
  }
}
\end{minted}

All events carry:

\begin{itemize}
  \item \texttt{eventType}: canonical event type (e.g.,
        \texttt{step/started}, \texttt{governance/approval\_required}).
  \item \texttt{timestamp}: RFC 3339 timestamp.
  \item \texttt{runId}: globally unique run identifier.
  \item \texttt{sequence}: monotonically increasing event sequence number
        (starts at 0).
  \item \texttt{payload}: event-specific data.
\end{itemize}

\paragraph{Reconnection protocol:}

If the WebSocket connection is lost, the client may reconnect and provide
\texttt{sinceSequence} to resume from the last received event:

\begin{minted}{json}
{
  "action": "subscribe",
  "runId": "run-abc123",
  "sinceSequence": 42  // Resume from event 43
}
\end{minted}

The server will replay all events with \texttt{sequence > 42}.

\paragraph{Stability guarantee:}

Within the \texttt{events/v2} contract:

\begin{itemize}
  \item Event envelope fields (\texttt{eventType}, \texttt{timestamp},
        \texttt{runId}, \texttt{sequence}) are stable.
  \item New event types may be added in minor versions.
  \item Existing event types will not change their \texttt{eventType} string.
  \item Payload schemas may add optional fields; clients MUST ignore unknown
        fields.
\end{itemize}

\subsection{Tool Invocation Contract (\texttt{tool/v2})}
\label{sec:tool-invocation-contract}

The tool invocation contract defines the stdio JSON-RPC protocol that tools
must implement to be invoked by gert.  All tools — whether built-in,
project-level, or extension-contributed — implement this contract.

\paragraph{Contract version:} \texttt{tool/v2}

\paragraph{Transport:} stdio JSON-RPC 2.0

\paragraph{Invocation request:}

\begin{minted}{json}
{
  "jsonrpc": "2.0",
  "id": "invocation-5",
  "method": "invoke",
  "params": {
    "action": "lookup",
    "inputs": {
      "id": "INC-12345"
    },
    "context": {
      "runId": "run-abc123",
      "stepId": "lookup_incident",
      "traceId": "trace-xyz789",
      "userId": "ops@acme.example",
      "mode": "real",
      "attempt": 1
    }
  }
}
\end{minted}

\paragraph{Success response:}

\begin{minted}{json}
{
  "jsonrpc": "2.0",
  "id": "invocation-5",
  "result": {
    "outputs": {
      "incident": {
        "id": "INC-12345",
        "severity": "high",
        "status": "open"
      }
    },
    "exitCode": 0,
    "telemetry": {
      "durationMs": 320,
      "startedAt": "2026-04-18T12:00:00Z",
      "finishedAt": "2026-04-18T12:00:00.320Z"
    }
  }
}
\end{minted}

\paragraph{Error response:}

\begin{minted}{json}
{
  "jsonrpc": "2.0",
  "id": "invocation-5",
  "error": {
    "code": -32050,
    "message": "incident not found",
    "data": {
      "category": "runtime",
      "retryable": false,
      "details": {
        "incidentId": "INC-12345",
        "statusCode": 404
      }
    }
  }
}
\end{minted}

\paragraph{Cancellation request:}

\begin{minted}{json}
{
  "jsonrpc": "2.0",
  "id": "cancel-5",
  "method": "cancel",
  "params": {
    "invocationId": "invocation-5",
    "reason": "run-cancelled"
  }
}
\end{minted}

\paragraph{MCP tool adapter:}

When an MCP server is used as a tool source (§5), the gert host wraps the MCP
\texttt{tools/call} protocol to satisfy the tool invocation contract.  The
adapter:

\begin{enumerate}
  \item Translates the \texttt{invoke} request to an MCP \texttt{tools/call}
        request.
  \item Forwards the \texttt{context} object as \texttt{\_gert\_context} in the
        MCP request (extension to MCP protocol).
  \item Translates the MCP response back to the tool invocation contract
        response.
  \item Flattens MCP \texttt{content} arrays to a single string for capture.
\end{enumerate}

\paragraph{Stability guarantee:}

Within the \texttt{tool/v2} contract:

\begin{itemize}
  \item Request/response envelope fields are stable.
  \item The \texttt{context} object fields are stable; new fields may be added.
  \item Error code ranges are stable (see Section~\ref{sec:error-codes}).
  \item Tools MUST ignore unknown fields in requests.
\end{itemize}

\subsection{Extension Handshake Contract (\texttt{extension/v2})}
\label{sec:extension-handshake-contract}

The extension handshake contract defines the initialization and capability
negotiation protocol between the gert host and an extension process.  It is
defined in §4 and summarized here for completeness.

\paragraph{Contract version:} \texttt{extension/v2}

\paragraph{Transport:} stdio JSON-RPC 2.0

\paragraph{Initialize request:}

\begin{minted}{json}
{
  "jsonrpc": "2.0",
  "id": "1",
  "method": "extension/initialize",
  "params": {
    "host": {
      "name": "gert",
      "version": "2.0.0",
      "apiVersion": "2.0"
    },
    "protocolVersion": "2.0",
    "grantedCapabilities": [
      "capability/tool-registration",
      "capability/network"
    ],
    "runContext": {
      "runId": "run-abc123",
      "mode": "real"
    }
  }
}
\end{minted}

\paragraph{Initialize response:}

\begin{minted}{json}
{
  "jsonrpc": "2.0",
  "id": "1",
  "result": {
    "extension": {
      "name": "acme.incident-pack",
      "version": "1.2.0"
    },
    "protocolVersion": "2.0",
    "acknowledgedCapabilities": [
      "capability/tool-registration",
      "capability/network"
    ]
  }
}
\end{minted}

\paragraph{Heartbeat (ping/pong):}

\begin{verbatim}
Request:
{
  "jsonrpc": "2.0",
  "id": "42",
  "method": "extension/ping",
  "params": {}
}

Response:
{
  "jsonrpc": "2.0",
  "id": "42",
  "result": {
    "status": "ok"
  }
}
\end{verbatim}

\paragraph{Stability guarantee:}

Within the \texttt{extension/v2} contract:

\begin{itemize}
  \item The \texttt{extension/initialize} request/response schema is stable.
  \item New capabilities may be added; extensions MUST ignore unknown
        capabilities in \texttt{grantedCapabilities}.
  \item The \texttt{extension/ping} method signature is stable.
  \item Extension error codes (§4) are stable.
\end{itemize}

% ─────────────────────────────────────────────────────────────────────────────
\section{Contract Versioning}
\label{sec:contract-versioning}
% ─────────────────────────────────────────────────────────────────────────────

Each adapter contract has a version declared in the initial handshake or HTTP
header.  The version format is \texttt{\{surface\}/\{major\}} (e.g.,
\texttt{exec/v2}, \texttt{events/v2}).  There is no minor version: any breaking
change bumps the major version.

\subsection{Version Declaration}

\paragraph{HTTP-based contracts (JSON-RPC, WebSocket):}

Clients MUST send an \texttt{X-Gert-Contract} header on the initial HTTP
request:

\begin{verbatim}
GET /ws HTTP/1.1
Host: gert.acme.example:8443
X-Gert-Contract: events/v2
\end{verbatim}

The server responds with:

\begin{verbatim}
HTTP/1.1 200 OK
X-Gert-Contract: events/v2
\end{verbatim}

If the client sends an unsupported contract version, the server responds with
HTTP 400 and error:

\begin{minted}{json}
{
  "error": {
    "code": 4001,
    "message": "unsupported contract version",
    "data": {
      "requestedVersion": "events/v3",
      "supportedVersions": ["events/v2"]
    }
  }
}
\end{minted}

\paragraph{stdio-based contracts (tool, extension):}

The contract version is declared in the first message exchanged (initialize or
handshake):

\begin{minted}{json}
{
  "protocolVersion": "2.0"
}
\end{minted}

If the protocol version is unsupported, the responder returns error code
\texttt{-32001} (protocol version not supported).

\subsection{Breaking vs. Non-Breaking Changes}

\paragraph{Non-breaking changes (minor version, no contract version bump):}

\begin{itemize}
  \item Adding optional fields to requests or responses.
  \item Adding new methods to the contract.
  \item Adding new error codes.
  \item Adding new event types.
\end{itemize}

Consumers MUST implement forward compatibility: ignore unknown fields, ignore
unknown methods, ignore unknown event types.

\paragraph{Breaking changes (major version bump):}

\begin{itemize}
  \item Removing or renaming required fields.
  \item Changing the semantics of existing methods.
  \item Removing methods.
  \item Changing error code meanings.
  \item Incompatible transport changes (e.g., switching from JSON-RPC to gRPC).
\end{itemize}

Breaking changes require a new contract version (e.g., \texttt{exec/v2} →
\texttt{exec/v3}).

\subsection{Deprecation Policy}

\begin{itemize}
  \item When a new major contract version is released, the previous version is
        supported for 12 months after the new version's GA date.
  \item During the deprecation period, both versions are supported
        simultaneously.
  \item After 12 months, the old version is removed.  Clients must upgrade.
\end{itemize}

Example timeline:

\begin{itemize}
  \item 2026-04-01: \texttt{exec/v2} released.
  \item 2027-06-01: \texttt{exec/v3} released; \texttt{exec/v2} enters
        deprecation.
  \item 2028-06-01: \texttt{exec/v2} removed; only \texttt{exec/v3} supported.
\end{itemize}

\subsection{Forward and Backward Compatibility}

\paragraph{Forward compatibility (clients):}

Clients MUST ignore unknown fields in responses.  This allows the server to add
optional fields without breaking existing clients.

Example:

\begin{minted}{bash}
// Server adds a new optional field in exec/v2
{
  "result": {
    "runId": "run-abc123",
    "status": "running",
    "newOptionalField": "value"  // Clients ignore this
  }
}
\end{minted}

\paragraph{Backward compatibility (servers):}

Servers MUST accept requests that omit optional fields added in later minor
versions.  Servers MUST NOT require new optional fields.

Example:

\begin{minted}{bash}
// Client sends an old request format (before a new optional field was added)
{
  "method": "exec/start",
  "params": {
    "runbookPath": "./runbook.yaml"
    // Missing new optional field "labels"
  }
}
// Server accepts and defaults the missing field
\end{minted}

% ─────────────────────────────────────────────────────────────────────────────
\section{Error Codes}
\label{sec:error-codes}
% ─────────────────────────────────────────────────────────────────────────────

All adapter contracts use a canonical error code registry.  Error codes are
integers organized by category.  Codes are stable across contract versions.

\begin{table}[ht]
\centering
\begin{tabular}{p{0.15\linewidth}p{0.60\linewidth}p{0.15\linewidth}}
\hline
\textbf{Code} & \textbf{Description} & \textbf{Retryable} \\
\hline
\multicolumn{3}{c}{\textbf{1xxx — Execution Errors}} \\
\hline
1001 & Run not found (invalid \texttt{runId}) & No \\
1002 & Step failed (runtime error in step execution) & Maybe \\
1003 & Run cancelled (user or timeout cancellation) & No \\
1004 & Run already exists (duplicate \texttt{runId}) & No \\
\hline
\multicolumn{3}{c}{\textbf{2xxx — Governance Errors}} \\
\hline
2001 & Command denied (not in allowlist or in denylist) & No \\
2002 & Approval timeout (no approval received in time) & No \\
2003 & Approval rejected (approver rejected the request) & No \\
2004 & Insufficient approvals (min approval count not reached) & No \\
\hline
\multicolumn{3}{c}{\textbf{3xxx — Schema Errors}} \\
\hline
3001 & Schema validation failed (runbook does not match schema) & No \\
3002 & Unknown step type (unrecognized \texttt{stepType}) & No \\
3003 & Import cycle detected (circular runbook dependencies) & No \\
\hline
\multicolumn{3}{c}{\textbf{4xxx — Integration Errors}} \\
\hline
4001 & Tool not found (tool name does not resolve) & No \\
4002 & Extension handshake failed (protocol error) & Maybe \\
4003 & Input resolution failed (provider error or missing input) & Maybe \\
4004 & Unsupported contract version & No \\
\hline
\multicolumn{3}{c}{\textbf{5xxx — Server Errors}} \\
\hline
5001 & Internal error (unexpected host failure) & Maybe \\
5002 & Resource exhausted (rate limit or quota exceeded) & Yes \\
5003 & Service unavailable (host shutting down or overloaded) & Yes \\
\hline
\end{tabular}
\caption{Canonical error code registry for adapter contracts.}
\label{tab:error-codes}
\end{table}

\paragraph{Error envelope structure:}

\begin{minted}{json}
{
  "error": {
    "code": 2001,
    "message": "command denied: not in allowlist",
    "data": {
      "category": "governance",
      "retryable": false,
      "details": {
        "command": "rm",
        "reason": "in_denylist"
      }
    }
  }
}
\end{minted}

\paragraph{Error code stability:}

\begin{itemize}
  \item Existing error codes will not change their meaning.
  \item New error codes may be added in minor versions.
  \item Clients MUST handle unknown error codes gracefully (fallback to generic
        error handling).
\end{itemize}

% ─────────────────────────────────────────────────────────────────────────────
\section{Contract Test Suite}
\label{sec:contract-test-suite}
% ─────────────────────────────────────────────────────────────────────────────

The gert repository includes a contract test suite that validates adapter
implementations against the canonical contracts.  Adapter authors SHOULD run
these tests to ensure compatibility.

\subsection{Test Suite Structure}

Contract tests are located in \texttt{testdata/contracts/}:

\begin{verbatim}
testdata/contracts/
  exec-v2/
    01-start-run.test.yaml
    02-cancel-run.test.yaml
    03-list-runs.test.yaml
  events-v2/
    01-subscribe.test.yaml
    02-reconnect.test.yaml
  tool-v2/
    01-invoke.test.yaml
    02-cancel.test.yaml
  extension-v2/
    01-initialize.test.yaml
    02-ping.test.yaml
\end{verbatim}

Each test case is a YAML file with:

\begin{itemize}
  \item \texttt{description}: human-readable test case description.
  \item \texttt{request}: the JSON-RPC request or WebSocket message.
  \item \texttt{expectedResponse}: the expected response (or error).
  \item \texttt{preconditions}: setup required before the test (e.g., "a run
        with id run-abc123 exists").
\end{itemize}

\subsection{Running Contract Tests}

\begin{minted}{bash}
$ gert contract test --adapter=vscode
Running contract tests for adapter: vscode
  ✓ exec-v2/01-start-run (passed)
  ✓ exec-v2/02-cancel-run (passed)
  ✗ events-v2/01-subscribe (failed: missing 'sequence' field)

Contract test results: 2 passed, 1 failed
\end{minted}

Adapter implementations SHOULD achieve 100\% pass rate before release.

\subsection{Contract Test CI Integration}

Contract tests are part of the CI gate for all adapter implementations:

\begin{itemize}
  \item The VS Code extension CI runs \texttt{gert contract test
        --adapter=vscode}.
  \item Third-party adapters are encouraged to use the same test suite.
  \item Failed contract tests block merges to the main branch.
\end{itemize}

% ─────────────────────────────────────────────────────────────────────────────
\section{Reference Implementations}
\label{sec:reference-implementations}
% ─────────────────────────────────────────────────────────────────────────────

\paragraph{VS Code extension.}

The official VS Code extension lives in its own repository at
\url{https://github.com/ormasoftchile/gert-vscode}. It is the reference implementation of
the HTTP transport surface exposed by \texttt{gert serve}: it manages a child
\texttt{gert serve} process, attaches a webview that loads
\texttt{/preview/} from the local server, and consumes the SSE event streams
(\texttt{/events}, \texttt{/runs/\{id\}/state}, \texttt{/runs/\{id\}/interactions}) plus
the REST/JSON-RPC endpoints. Adapter authors should study its implementation for best
practices. The extension is packaged with \texttt{vsce} and ships a \texttt{.vsix}
artifact from CI.

\paragraph{Built-in tools.}

The built-in tools (\texttt{ext/builtin/tools/} in the repository) are the
reference implementation of the tool invocation contract (\texttt{tool/v2}).
They demonstrate how to handle \texttt{invoke}, \texttt{cancel}, and error
responses.

\paragraph{MCP adapter.}

The MCP tool adapter (\texttt{ext/toolruntime/mcp\_adapter.go}) demonstrates
how to wrap an external protocol (MCP) to satisfy the tool invocation contract.

% ─────────────────────────────────────────────────────────────────────────────
\section{Contract Stability Guarantees}
% ─────────────────────────────────────────────────────────────────────────────

Gert v2 commits to the following stability guarantees for all adapter
contracts:

\begin{enumerate}
  \item \textbf{Semantic versioning}: major version bumps indicate breaking
        changes; minor/patch versions are backward compatible.
  \item \textbf{12-month deprecation period}: old contract versions are
        supported for 12 months after a new major version is released.
  \item \textbf{Additive changes only}: within a major version, only
        non-breaking changes (adding optional fields, new methods, new event
        types) are permitted.
  \item \textbf{Error code stability}: error codes do not change meaning within
        a major version.
  \item \textbf{Forward compatibility}: clients must ignore unknown fields;
        servers must accept omitted optional fields.
\end{enumerate}

These guarantees are designed to minimize adapter churn and allow external
systems to depend on gert contracts with confidence.

% ─────────────────────────────────────────────────────────────────────────────
\section{Contract Documentation and Tooling}
% ─────────────────────────────────────────────────────────────────────────────

\subsection{OpenAPI Specifications}

All HTTP-based contracts (JSON-RPC over HTTP, WebSocket event stream) are
documented as OpenAPI 3.1 specifications~\cite{openapi-spec} in
\texttt{specs/contracts/}:

\begin{itemize}
  \item \texttt{specs/contracts/exec-v2.openapi.yaml}
  \item \texttt{specs/contracts/events-v2.openapi.yaml}
\end{itemize}

These specifications are the source of truth for request/response schemas and
are used to generate client libraries.

\subsection{JSON Schema Validation}

All JSON-RPC request/response schemas are defined as JSON Schema (Draft
2020-12)~\cite{jsonschema-spec} in \texttt{specs/contracts/schemas/}.  These
schemas are used for:

\begin{itemize}
  \item Runtime validation of requests/responses.
  \item Contract test validation.
  \item Client library generation (TypeScript, Python, Go).
\end{itemize}

\subsection{Client Library Generation}

The gert project provides code generation tools to generate client libraries
from the OpenAPI and JSON Schema definitions:

\begin{minted}{bash}
$ gert codegen client --language=typescript --output=./clients/ts/
Generated TypeScript client for exec/v2 and events/v2
Files written to ./clients/ts/
\end{minted}

Official client libraries are provided for:

\begin{itemize}
  \item TypeScript (for VS Code and web adapters)
  \item Python (for CI/CD integrations)
  \item Go (for extension authors)
\end{itemize}
